\section{A seven-polynomial word with a complete list of size seven}
\label{app:seven-cubic-bank}

This construction gives an exact list, rather than only seven exhibited
nearby polynomials.  Its list size is fixed; no assertion about an unbounded
list or about exceptional labels on a received affine line follows from it.

\begin{theorem}\label{thm:seven-cubic-bank}
For every integer $m\geq1$, there are a number field $K_m$, a set
$\Omega_m\subset K_m$ of $14m$ distinct points, and a word on $\Omega_m$
whose complete list of polynomials of degree at most $3m$ with at least
$7m$ agreements has exactly seven elements.  Each of those seven
polynomials has exactly $7m$ agreements.  For each fixed $m$, the same
statement holds over $\mathbb F_p$ for arbitrarily large primes $p$.
In particular, $m=2$ gives length $28$, dimension $7$, agreement $14$,
and complete list size seven, at exactly quarter rate.
\end{theorem}

\begin{proof}
Let $\theta$ be a root of
\[
 f(T)=T^3+2T^2-T-1,
\]
and work first over $K=\mathbb Q(\theta)$.  The polynomial $f$ is
irreducible: its only possible rational roots are $1,-1$, and neither
is a root.  Put
\[
 u=\theta^2,\quad v=\theta^2+\theta-1,\quad z=\theta^2+\theta,
 \qquad j=1,\quad k=3\theta^2+4\theta-5,\quad \ell=k-1.
\]
We initially evaluate binary cubics on the projective line; at infinity,
the value of an affine polynomial is its coefficient of $X^3$.
Set $P_1=0$ and define
\[
\begin{array}{lll}
 Q_2=(X-\theta)(X-z),&Q_3=(X-v)(X-z),&Q_4=(X-u)(X-v),\\
 Q_5=(X-u)(X-z),&Q_6=(X-u)(X-\theta),&Q_7=(X-v)(X-\theta).
\end{array}
\]
The other six cubics are
\begin{align*}
 P_i&=Q_i\left[\left(\frac{k}{Q_i(1)}-\frac{j}{Q_i(0)}\right)X
                  +\frac{j}{Q_i(0)}\right] &&(i=2,4),\\
 P_i&=Q_i\left[\ell(X-1)+\frac{k}{Q_i(1)}\right] &&(i=3,6),\\
 P_i&=Q_i\left[\ell X+\frac{j}{Q_i(0)}\right] &&(i=5,7).
\end{align*}
All these denominators are nonzero.

Here is the entire incidence certificate.  A string such as $123$
denotes the set $\{1,2,3\}$.  At each coordinate the received value is
the common value of the indicated polynomials.
\[
\begin{array}{c|c@{\qquad}c|c}
 \text{coordinate}&\text{matching indices}&
 \text{coordinate}&\text{matching indices}\\ \hline
 \infty&3567&1-\theta^2&123\\
 0&2457&\theta^2+2\theta-1&145\\
 1&2346&\theta^2+2\theta+1&167\\
 u&1456&\theta^2-1&246\\
 v&1347&-\theta^2-2\theta+3&257\\
 \theta&1267&3\theta^2-1&347\\
 z&1235&(3\theta^2+4\theta-1)/7&356
\end{array}
\]
The fourteen coordinates are distinct.  Substitution in the displayed
formulas, followed by reduction using $f(\theta)=0$, verifies every
equality in the table and verifies that there are no additional matches.
In particular the seven cubics are distinct.  Each index occurs in four
of the left-hand quadruples and three of the right-hand triples, hence
each polynomial has exactly seven agreements.  This certificate consists
of identities in the explicit cubic field; it requires no genericity
assumption or geometric classification of other configurations.

We next show that the list is complete, even over an algebraic closure.
If an eighth distinct cubic had at least seven agreements, and $s_x$
were the number of matching polynomials among these eight at coordinate
$x$, then
\[
 \sum_x s_x\geq56,
 \qquad
 \sum_x\binom{s_x}{2}\leq3\binom82=84.
\]
Convexity gives the reverse bound $\sum_x\binom{s_x}{2}\geq84$,
with equality only if all fourteen $s_x$ equal four.  Thus the eighth
cubic would have to match exactly the seven triple coordinates in the
right-hand column of the table.
Let those coordinates be $a_1,\ldots,a_7$, with received values
$r_1,\ldots,r_7$.  The leading coefficient of their degree-at-most-six
Lagrange interpolant $R$ is
\[
 [X^6]R
 =\sum_{i=1}^7\frac{r_i}{\prod_{h\ne i}(a_i-a_h)}
 =\frac{42\theta^2+81\theta-30}{8}\ne0.
\]
The equality is another direct reduction modulo $f$; nonvanishing follows
from its irreducibility.  Consequently no cubic interpolates those seven
values.  An eighth candidate, and hence any larger list, is impossible.

For an affine domain, change coordinates by
\[
 t=\frac1{X-2},\qquad H_i(t)=t^3P_i(2+1/t).
\]
None of the fourteen coordinates equals $2$: every nonconstant finite
coordinate in the table has degree at most two in $\theta$.
This change therefore gives fourteen distinct finite coordinates, with
the old infinity mapped to zero.  Scale the received value at a finite
old coordinate by $t^3$, and retain the old infinity value at $t=0$.
The $H_i$ have degree at most three and all agreements are preserved.
Call this affine domain $\Omega$ and its word $r$.

For the general construction choose a degree-$m$ polynomial $\psi$
such that the fourteen fibers over $\Omega$ are disjoint and each has
$m$ distinct points.  Such a choice exists over a finite extension of
$K$: for example $\psi(U)=U^m+c$ with $c\notin\Omega$ suffices in
characteristic zero.  On $\Omega_m=\psi^{-1}(\Omega)$ use the word
$r\circ\psi$ and the seven polynomials $H_i\circ\psi$.
Their matching sets have sizes $7m$.
If an eighth polynomial of degree at most $3m$ had at least $7m$
agreements, the preceding pair count, multiplied by $m$, would again
force it to match every point above the seven triple coordinates and
none above the quadruple coordinates.
Let $\widetilde R$ be the interpolant of the seven triple values on the
affine domain.  Its degree lies between four and six: a cubic interpolant
would contradict the projective argument just proved.
The difference between the alleged eighth polynomial and
$\widetilde R\circ\psi$ has at least $7m$ distinct roots and degree at
most $6m$, so it is zero.  This is impossible because
$\deg(\widetilde R\circ\psi)>3m$.  Thus the pulled-back list is also
exactly seven.

Finally all coefficients and fiber points lie in a number field.
Exclude the finitely many primes dividing their denominators, the
nonzero coordinate differences, the nonzero incidence differences,
and the relevant leading interpolation coefficients.  At primes
splitting completely in a Galois closure, reduction then preserves the
entire construction and the complete-list proof.  There are arbitrarily
large such primes.  This proves the prime-field assertion.
\end{proof}

For a small explicit instance with $m=1$, take $p=97$.  In the displayed
order, the domain and word are
\[
\begin{split}
 \Omega={}&(0,48,96,91,79,11,41,51,52,55,13,42,78,21),\\
 r={}&(74,12,22,0,0,0,0,0,0,0,72,50,91,75).
\end{split}
\]
The seven coefficient vectors, in increasing degree order, are
\[
\begin{gathered}
 (0,0,0,0),\ (47,12,72,85),\ (74,90,75,37),\
 (40,96,82,4),\\
 (74,26,3,10),\ (74,46,10,16),\ (74,27,35,70).
\end{gathered}
\]
They arise from $\theta=55$ in the construction.  The seven-point
interpolation obstruction above reduces to $5\ne0$ modulo $97$, so this
example also has complete list size exactly seven.
